%!TEX root = ../main.tex
\section{Related Work}
\label{sec:related-work}

\textbf{Policy gradient methods.}
In addition to the algorithms mentioned in Section~\ref{sec:intro}, representative policy-gradient variants include Natural Policy Gradient (NPG) \cite{kakade01neurips-npg}, Deep Deterministic Policy Gradient (DDPG) \cite{lillicrap15arxiv-ddpg}, and Soft Actor--Critic (SAC) \cite{haarnoja18icml-sac}. NPG preconditions the gradient using the Fisher information matrix, but it still estimates gradients using REINFORCE-style likelihood-ratio estimators; therefore, it falls within the scope of our analysis. In contrast, DDPG and SAC are actor--critic methods and use critic-based gradient surrogates.

For policy-gradient methods on LQR/LQG, \cite{fazel18icml-lqrglobal} study infinite-horizon LQR with fixed Gaussian exploration covariance from an optimization perspective, showing that the objective is gradient dominated and that policy gradient converges to the global optimum. \cite{furieri20l4dc-finitelqr} extends these guarantees to a class of finite-horizon LQR problems. \cite{malik20jmlr-stochasticlqroptimization,mohammadi21tac-stochasticlqroptimization} further consider stochastic estimation settings and establish convergence and sample complexity results. \cite{yang19neurips-actorcriticlqr} analyzes the convergence of actor--critic algorithms on LQR. For partially observed LQG, \cite{tang23mp-lqglandscape} characterizes the optimization landscape, including connectivity properties of stable controllers and the structure of stationary points. Beyond LQR/LQG, \cite{mei20icml-globalconvergence,agarwal21jmlr-tabularglobalconveregnce} study global convergence in tabular settings and derive error bounds in terms of distribution shift/transfer.

More closely related to our focus on estimator variability, \cite{zhao11neurips-pgvariance} shows that parameter-based exploration can yield lower-variance policy-gradient estimates than the vanilla REINFORCE estimator; both admit $O(\sigma^{-2})$ upper and lower bounds, where $\sigma$ denotes the standard deviation of a Gaussian policy. Their analysis assumes bounded states and rewards and does not characterize an exact NSR expression. \cite{preiss19arxiv-lqgvariance} derives variance upper bounds for finite-horizon LQG and studies contributing factors empirically, but focuses on stable systems and does not make the scaling with the initial-state distribution or policy covariance explicit. We complement these works by providing exact finite-horizon variance expressions and by extending the analysis to polynomial systems and to generic nonlinear systems.

\textbf{Convergence dynamics.}
\cite{ilyas18arxiv-policycollapse} empirically observes that the cosine similarity between the estimated and true gradients can deteriorate as optimization progresses; we interpret this phenomenon through our NSR analysis. \cite{dohare23ewrl-policycollapse} further reports a ``policy collapse'' phenomenon and relates it to large Adam momentum. A complementary line of work studies how baseline choices affect optimization beyond their impact on variance: \cite{chung21icml-committalbaselines} shows that even when the variance is held fixed, different baseline choices can lead to markedly different convergence behavior, and \cite{mei22neurips-committalbaselines} establishes convergence results in bandit settings while showing that the variance near the optimal policy can be unbounded for any choice of baseline.


% The study of policy gradient can be divided into two groups: the theoretical analysis of policy gradient's variance and convergence of REINFORCE algorithm; and the practical algorithms that improve the stability of policy gradient methods, such as TRPO \cite{schulman15icml-trpo} and PPO \cite{schulman17arxiv-ppo}.

% Trust Region Policy Optimization (TRPO) \cite{schulman15icml-trpo} firstly introduce performance difference lemma, which indicates that the performance of a new policy can be reformulated as the advantage of the new policy but under current policy's state distribution. Based on this lemma, TRPO propose to optimize a surrogate objective with a constraint on the KL divergence between the new and old policy. This made the new state distribution close to the old one, thus ensuring monotonic improvement of the policy. We call these method "curvature analysis" since they all improving the curvature on the decent direction. 

% Fazel et al. \cite{fazel18icml-lqrglobal} first study policy gradient methods on Linear Quadratic Regulator (LQR) from an optimization perspective. They show that the LQR objective is gradient dominated and can converge to the global optimum. We follow their setup and derive the exact analytical gradient for infinite-horizon discounted LQR problem with Gaussian policies, but further consider addtional factors like finite horizon, gaussian policy, discounted reward. Preiss et al. \cite{preiss19arxiv-lqgvariance} study the convergence of policy gradient methods on Linear Quadratic Gaussian (LQG) problem, which is more close to our setting. They show that the variance increase polynomially with respect to the horizon, and goes to infty if spectral radius is approximating 1. We complement their work by providing exact variance expression with finite horizon, whatever the system is stable or not.

% While this zero-order method is easy to implement, it often suffers from high variance. An altenative method is \textit{Pathwise} gradient estimator \cite{mohamed20jmlr-montecarlo,schulman15neurips-stocompgraph}, which leverages differentiable simulator to backpropagate the reward signal through the trajectory. This method can significantly reduce variance, but requires access to environment model.

% \begin{takeawaybox}{Path from deterministic optimization to variance analysis}
%     This report originated from replacing the expectation with a finite-sample average and studying the optimization properties of the resulting deterministic objective $J_{\text{opt}} = \sum_i R(\tau(s_0^{(i)}))$. Because both the dynamics and the policy are deterministic, the entire trajectory $\tau$ is determined by the initial state $s_0^{(i)}$. We employed the bundle method, a state-of-the-art nonconvex optimization technique, to solve the empirical objective, but found convergence brittle and gradients very noisy. A brief result is shown in Figure.~\ref{fig:bundle_landscape}.
%     \begin{figure}[H]
%         \centering
%         \includegraphics[width=1.0\textwidth]{figure/bundle_landscape.pdf}
%         \caption{After running the bundle method on the finite-sample objective, the landscape appears noisy and the gradient brittle. As the number of samples increases, the landscape approaches the true expectation, while the gradient remains noisy.}
%         \label{fig:bundle_landscape}
%     \end{figure}
%     We then compared different policy-gradient estimators on the Acrobot problem and observed that estimates near the optimal policy are highly inaccurate. This motivates a deeper study of the variance behavior of policy gradients.
% \end{takeawaybox}
    